\section{Erweiterungs- und Verbesserungspotenziale}
Die entwickelte WordClock bietet verschiedene Möglichkeiten zur technischen Weiterentwicklung.  
Durch die modulare Architektur können zusätzliche Funktionen vergleichsweise einfach integriert werden.  

Eine mögliche Erweiterung betrifft die Optimierung der Energieeffizienz.  
Hierzu könnten adaptive Helligkeitsregelungen, Energiesparmodi oder effizientere Verfahren zur \gls{led}-Ansteuerung implementiert werden.  
Ebenso wäre der Einsatz zusätzlicher Sensorik zur automatischen Anpassung an Umgebungsbedingungen möglich.  

Auch die Netzwerkfunktionen können erweitert werden.  
Denkbar sind beispielsweise Smart-Home-Schnittstellen, Cloud-Anbindungen oder die Integration weiterer Online-Dienste.  
Zusätzlich könnte eine mobile Anwendung zur erweiterten Steuerung und Konfiguration entwickelt werden.  

Im Bereich der Anzeige bestehen weitere Gestaltungsmöglichkeiten.  
Hierzu zählen zusätzliche Animationen, individuell anpassbare Designs oder alternative Darstellungsmodi.  
Ebenso wäre eine größere oder höher auflösende Anzeige denkbar.  

Darüber hinaus könnte die mechanische Konstruktion weiter optimiert werden.  
Mögliche Verbesserungen betreffen die Montagefreundlichkeit, die Wärmeabfuhr sowie die Fertigungsqualität einzelner Bauteile.  
Auch alternative Fertigungsverfahren oder Materialien könnten untersucht werden.  

Für zukünftige Entwicklungen wäre zusätzlich eine umfangreiche Langzeituntersuchung sinnvoll.  
Dabei könnten Aspekte wie Zuverlässigkeit, thermisches Verhalten, Energieverbrauch und Dauerstabilität systematisch analysiert werden.  

Zusammenfassend bietet die entwickelte WordClock eine geeignete Grundlage für weiterführende technische Entwicklungen.  
Das System zeigt das Potenzial moderner, vernetzter und interaktiver Anzeige- und Informationssysteme innerhalb elektrotechnischer Anwendungen.